% Automatic Section Continuation Markers
% ============================================================================
%
% PURPOSE:
%   Automatically insert continuation markers (e.g., "(Continued)") when
%   nested sections close and content resumes at a parent section level.
%
% REQUIREMENTS:
%   - This file provides the `asection` environment so no external package is needed
%   - LaTeX3 (expl3) is required for the peek logic that checks following content
%
% USAGE:
%   1. Include this file in your preamble:
%      % Automatic Section Continuation Markers
% ============================================================================
%
% PURPOSE:
%   Automatically insert continuation markers (e.g., "(Continued)") when
%   nested sections close and content resumes at a parent section level.
%
% REQUIREMENTS:
%   - This file provides the `asection` environment so no external package is needed
%   - LaTeX3 (expl3) is required for the peek logic that checks following content
%
% USAGE:
%   1. Include this file in your preamble:
%      % Automatic Section Continuation Markers
% ============================================================================
%
% PURPOSE:
%   Automatically insert continuation markers (e.g., "(Continued)") when
%   nested sections close and content resumes at a parent section level.
%
% REQUIREMENTS:
%   - This file provides the `asection` environment so no external package is needed
%   - LaTeX3 (expl3) is required for the peek logic that checks following content
%
% USAGE:
%   1. Include this file in your preamble:
%      % Automatic Section Continuation Markers
% ============================================================================
%
% PURPOSE:
%   Automatically insert continuation markers (e.g., "(Continued)") when
%   nested sections close and content resumes at a parent section level.
%
% REQUIREMENTS:
%   - This file provides the `asection` environment so no external package is needed
%   - LaTeX3 (expl3) is required for the peek logic that checks following content
%
% USAGE:
%   1. Include this file in your preamble:
%      \input{preamble/continuation-markers}
%
%   2. Use the supplied `asection` environment for nested sections:
%      \begin{asection}{Parent Section}
%      Parent content
%
%        \begin{asection}{Subsection}
%        Subsection content
%        \end{asection}
%
%        % Automatic "(Continued)" marker appears here (only if content follows)
%
%      More parent content
%      \end{asection}
%
% CUSTOMISATION:
%   Change marker text using \setcontinuationmarker:
%     \setcontinuationmarker{(Resumed)}
%
%   Or redefine \continuationmarker directly:
%     \renewcommand{\continuationmarker}{\subsection*{\textit{(Discussion Continues)}}}
%
% IMPLEMENTATION:
%   Uses LaTeX3 expl3 peek functions to detect content:
%   - asectiondepth: Current nesting level (0 = no active sections)
%   - maxasectiondepth: Deepest level encountered (detects level transitions)
%   - Peek-ahead after \end{asection} to check if another \begin{asection} follows
%
%   When asection closes and depth < maxdepth, peek ahead (ignoring spaces/newlines).
%   Only insert marker if the next non-space token is NOT \begin or \end.
%
% VERSION: 2.0.3
% DATE: 2025-11-13
% AUTHOR: Implementation based on feature specification 006-latex-section-continuation
% ============================================================================

\ExplSyntaxOn

% Counter Definitions
% ----------------------------------------------------------------------------
\newcounter{asectiondepth}               % Current nesting level (0 = no active sections)
\newcounter{maxasectiondepth}            % Maximum depth encountered (for detecting level transitions)

% Boolean flag to track if we need to check for continuation
% Use global (g) variable so it persists across environment boundaries
\bool_new:N \g_contmark_should_check_bool
\bool_set_false:N \g_contmark_should_check_bool

% Continuation Marker Commands
% ----------------------------------------------------------------------------

% \continuationmarker - Default continuation marker (unnumbered subsection)
% Users can customise by redefining this command
\cs_new_protected:Npn \continuationmarker { \subsection*{(Continued)} }

% \setcontinuationmarker{text} - Convenience command for customising marker text
\cs_new_protected:Npn \setcontinuationmarker #1 {
  \cs_set_protected:Npn \continuationmarker { \subsection*{#1} }
}

% Content Detection Logic
% ----------------------------------------------------------------------------

% Check if marker should be inserted by peeking ahead
% This is called when returning to a shallower nesting level
\cs_new_protected:Npn \__contmark_check_and_insert:
  {
    \bool_if:NT \g_contmark_should_check_bool
      {
        \bool_gset_false:N \g_contmark_should_check_bool
        \peek_meaning_ignore_spaces:NTF \begin
          {
            % Next command is \begin (likely another asection) - don't insert marker
          }
          {
            \peek_meaning_ignore_spaces:NTF \end
              {
                % Parent section closing - don't insert marker
              }
              {
                % Actual content follows - insert marker
                \continuationmarker
              }
          }
      }
  }

% Section dispatch logic
% ----------------------------------------------------------------------------
\cs_new_protected:Npn \__contmark_dispatch_section:nn #1#2
  {
    \int_case:nnF {#1}
      {
        {1}{\section{#2}}
        {2}{\subsection{#2}}
        {3}{\subsubsection{#2}}
        {4}{\paragraph{#2}}
        {5}{\subparagraph{#2}}
        {6}{\subsubparagraph{#2}}
      }
      {\paragraph{#2}}
  }

% Nesting management
% ----------------------------------------------------------------------------
\cs_new_protected:Npn \__contmark_begin_section:n #1
  {
    \stepcounter{asectiondepth}
    \int_compare:nNnT { \value{asectiondepth} } > { \value{maxasectiondepth} }
      {
        \setcounter{maxasectiondepth}{\value{asectiondepth}}
      }
    \__contmark_dispatch_section:nn { \value{asectiondepth} } {#1}
  }

\cs_new_protected:Npn \__contmark_end_section:
  {
    \addtocounter{asectiondepth}{-1}
    \int_compare:nNnT { \value{asectiondepth} } < { \value{maxasectiondepth} }
      {
        \int_compare:nNnT { \value{asectiondepth} } > { 0 }
          {
            \bool_gset_true:N \g_contmark_should_check_bool
            \setcounter{maxasectiondepth}{\value{asectiondepth}}
          }
      }
    \__contmark_check_and_insert:
  }

\newenvironment{asection}[1]
  {
    \__contmark_begin_section:n {#1}
  }
  {
    \__contmark_end_section:
  }
\ExplSyntaxOff

%
%   2. Use the supplied `asection` environment for nested sections:
%      \begin{asection}{Parent Section}
%      Parent content
%
%        \begin{asection}{Subsection}
%        Subsection content
%        \end{asection}
%
%        % Automatic "(Continued)" marker appears here (only if content follows)
%
%      More parent content
%      \end{asection}
%
% CUSTOMISATION:
%   Change marker text using \setcontinuationmarker:
%     \setcontinuationmarker{(Resumed)}
%
%   Or redefine \continuationmarker directly:
%     \renewcommand{\continuationmarker}{\subsection*{\textit{(Discussion Continues)}}}
%
% IMPLEMENTATION:
%   Uses LaTeX3 expl3 peek functions to detect content:
%   - asectiondepth: Current nesting level (0 = no active sections)
%   - maxasectiondepth: Deepest level encountered (detects level transitions)
%   - Peek-ahead after \end{asection} to check if another \begin{asection} follows
%
%   When asection closes and depth < maxdepth, peek ahead (ignoring spaces/newlines).
%   Only insert marker if the next non-space token is NOT \begin or \end.
%
% VERSION: 2.0.3
% DATE: 2025-11-13
% AUTHOR: Implementation based on feature specification 006-latex-section-continuation
% ============================================================================

\ExplSyntaxOn

% Counter Definitions
% ----------------------------------------------------------------------------
\newcounter{asectiondepth}               % Current nesting level (0 = no active sections)
\newcounter{maxasectiondepth}            % Maximum depth encountered (for detecting level transitions)

% Boolean flag to track if we need to check for continuation
% Use global (g) variable so it persists across environment boundaries
\bool_new:N \g_contmark_should_check_bool
\bool_set_false:N \g_contmark_should_check_bool

% Continuation Marker Commands
% ----------------------------------------------------------------------------

% \continuationmarker - Default continuation marker (unnumbered subsection)
% Users can customise by redefining this command
\cs_new_protected:Npn \continuationmarker { \subsection*{(Continued)} }

% \setcontinuationmarker{text} - Convenience command for customising marker text
\cs_new_protected:Npn \setcontinuationmarker #1 {
  \cs_set_protected:Npn \continuationmarker { \subsection*{#1} }
}

% Content Detection Logic
% ----------------------------------------------------------------------------

% Check if marker should be inserted by peeking ahead
% This is called when returning to a shallower nesting level
\cs_new_protected:Npn \__contmark_check_and_insert:
  {
    \bool_if:NT \g_contmark_should_check_bool
      {
        \bool_gset_false:N \g_contmark_should_check_bool
        \peek_meaning_ignore_spaces:NTF \begin
          {
            % Next command is \begin (likely another asection) - don't insert marker
          }
          {
            \peek_meaning_ignore_spaces:NTF \end
              {
                % Parent section closing - don't insert marker
              }
              {
                % Actual content follows - insert marker
                \continuationmarker
              }
          }
      }
  }

% Section dispatch logic
% ----------------------------------------------------------------------------
\cs_new_protected:Npn \__contmark_dispatch_section:nn #1#2
  {
    \int_case:nnF {#1}
      {
        {1}{\section{#2}}
        {2}{\subsection{#2}}
        {3}{\subsubsection{#2}}
        {4}{\paragraph{#2}}
        {5}{\subparagraph{#2}}
        {6}{\subsubparagraph{#2}}
      }
      {\paragraph{#2}}
  }

% Nesting management
% ----------------------------------------------------------------------------
\cs_new_protected:Npn \__contmark_begin_section:n #1
  {
    \stepcounter{asectiondepth}
    \int_compare:nNnT { \value{asectiondepth} } > { \value{maxasectiondepth} }
      {
        \setcounter{maxasectiondepth}{\value{asectiondepth}}
      }
    \__contmark_dispatch_section:nn { \value{asectiondepth} } {#1}
  }

\cs_new_protected:Npn \__contmark_end_section:
  {
    \addtocounter{asectiondepth}{-1}
    \int_compare:nNnT { \value{asectiondepth} } < { \value{maxasectiondepth} }
      {
        \int_compare:nNnT { \value{asectiondepth} } > { 0 }
          {
            \bool_gset_true:N \g_contmark_should_check_bool
            \setcounter{maxasectiondepth}{\value{asectiondepth}}
          }
      }
    \__contmark_check_and_insert:
  }

\newenvironment{asection}[1]
  {
    \__contmark_begin_section:n {#1}
  }
  {
    \__contmark_end_section:
  }
\ExplSyntaxOff

%
%   2. Use the supplied `asection` environment for nested sections:
%      \begin{asection}{Parent Section}
%      Parent content
%
%        \begin{asection}{Subsection}
%        Subsection content
%        \end{asection}
%
%        % Automatic "(Continued)" marker appears here (only if content follows)
%
%      More parent content
%      \end{asection}
%
% CUSTOMISATION:
%   Change marker text using \setcontinuationmarker:
%     \setcontinuationmarker{(Resumed)}
%
%   Or redefine \continuationmarker directly:
%     \renewcommand{\continuationmarker}{\subsection*{\textit{(Discussion Continues)}}}
%
% IMPLEMENTATION:
%   Uses LaTeX3 expl3 peek functions to detect content:
%   - asectiondepth: Current nesting level (0 = no active sections)
%   - maxasectiondepth: Deepest level encountered (detects level transitions)
%   - Peek-ahead after \end{asection} to check if another \begin{asection} follows
%
%   When asection closes and depth < maxdepth, peek ahead (ignoring spaces/newlines).
%   Only insert marker if the next non-space token is NOT \begin or \end.
%
% VERSION: 2.0.3
% DATE: 2025-11-13
% AUTHOR: Implementation based on feature specification 006-latex-section-continuation
% ============================================================================

\ExplSyntaxOn

% Counter Definitions
% ----------------------------------------------------------------------------
\newcounter{asectiondepth}               % Current nesting level (0 = no active sections)
\newcounter{maxasectiondepth}            % Maximum depth encountered (for detecting level transitions)

% Boolean flag to track if we need to check for continuation
% Use global (g) variable so it persists across environment boundaries
\bool_new:N \g_contmark_should_check_bool
\bool_set_false:N \g_contmark_should_check_bool

% Continuation Marker Commands
% ----------------------------------------------------------------------------

% \continuationmarker - Default continuation marker (unnumbered subsection)
% Users can customise by redefining this command
\cs_new_protected:Npn \continuationmarker { \subsection*{(Continued)} }

% \setcontinuationmarker{text} - Convenience command for customising marker text
\cs_new_protected:Npn \setcontinuationmarker #1 {
  \cs_set_protected:Npn \continuationmarker { \subsection*{#1} }
}

% Content Detection Logic
% ----------------------------------------------------------------------------

% Check if marker should be inserted by peeking ahead
% This is called when returning to a shallower nesting level
\cs_new_protected:Npn \__contmark_check_and_insert:
  {
    \bool_if:NT \g_contmark_should_check_bool
      {
        \bool_gset_false:N \g_contmark_should_check_bool
        \peek_meaning_ignore_spaces:NTF \begin
          {
            % Next command is \begin (likely another asection) - don't insert marker
          }
          {
            \peek_meaning_ignore_spaces:NTF \end
              {
                % Parent section closing - don't insert marker
              }
              {
                % Actual content follows - insert marker
                \continuationmarker
              }
          }
      }
  }

% Section dispatch logic
% ----------------------------------------------------------------------------
\cs_new_protected:Npn \__contmark_dispatch_section:nn #1#2
  {
    \int_case:nnF {#1}
      {
        {1}{\section{#2}}
        {2}{\subsection{#2}}
        {3}{\subsubsection{#2}}
        {4}{\paragraph{#2}}
        {5}{\subparagraph{#2}}
        {6}{\subsubparagraph{#2}}
      }
      {\paragraph{#2}}
  }

% Nesting management
% ----------------------------------------------------------------------------
\cs_new_protected:Npn \__contmark_begin_section:n #1
  {
    \stepcounter{asectiondepth}
    \int_compare:nNnT { \value{asectiondepth} } > { \value{maxasectiondepth} }
      {
        \setcounter{maxasectiondepth}{\value{asectiondepth}}
      }
    \__contmark_dispatch_section:nn { \value{asectiondepth} } {#1}
  }

\cs_new_protected:Npn \__contmark_end_section:
  {
    \addtocounter{asectiondepth}{-1}
    \int_compare:nNnT { \value{asectiondepth} } < { \value{maxasectiondepth} }
      {
        \int_compare:nNnT { \value{asectiondepth} } > { 0 }
          {
            \bool_gset_true:N \g_contmark_should_check_bool
            \setcounter{maxasectiondepth}{\value{asectiondepth}}
          }
      }
    \__contmark_check_and_insert:
  }

\newenvironment{asection}[1]
  {
    \__contmark_begin_section:n {#1}
  }
  {
    \__contmark_end_section:
  }
\ExplSyntaxOff

%
%   2. Use the supplied `asection` environment for nested sections:
%      \begin{asection}{Parent Section}
%      Parent content
%
%        \begin{asection}{Subsection}
%        Subsection content
%        \end{asection}
%
%        % Automatic "(Continued)" marker appears here (only if content follows)
%
%      More parent content
%      \end{asection}
%
% CUSTOMISATION:
%   Change marker text using \setcontinuationmarker:
%     \setcontinuationmarker{(Resumed)}
%
%   Or redefine \continuationmarker directly:
%     \renewcommand{\continuationmarker}{\subsection*{\textit{(Discussion Continues)}}}
%
% IMPLEMENTATION:
%   Uses LaTeX3 expl3 peek functions to detect content:
%   - asectiondepth: Current nesting level (0 = no active sections)
%   - maxasectiondepth: Deepest level encountered (detects level transitions)
%   - Peek-ahead after \end{asection} to check if another \begin{asection} follows
%
%   When asection closes and depth < maxdepth, peek ahead (ignoring spaces/newlines).
%   Only insert marker if the next non-space token is NOT \begin or \end.
%
% VERSION: 2.0.3
% DATE: 2025-11-13
% AUTHOR: Implementation based on feature specification 006-latex-section-continuation
% ============================================================================

\ExplSyntaxOn

% Counter Definitions
% ----------------------------------------------------------------------------
\newcounter{asectiondepth}               % Current nesting level (0 = no active sections)
\newcounter{maxasectiondepth}            % Maximum depth encountered (for detecting level transitions)

% Boolean flag to track if we need to check for continuation
% Use global (g) variable so it persists across environment boundaries
\bool_new:N \g_contmark_should_check_bool
\bool_set_false:N \g_contmark_should_check_bool

% Continuation Marker Commands
% ----------------------------------------------------------------------------

% \continuationmarker - Default continuation marker (unnumbered subsection)
% Users can customise by redefining this command
\cs_new_protected:Npn \continuationmarker { \subsection*{(Continued)} }

% \setcontinuationmarker{text} - Convenience command for customising marker text
\cs_new_protected:Npn \setcontinuationmarker #1 {
  \cs_set_protected:Npn \continuationmarker { \subsection*{#1} }
}

% Content Detection Logic
% ----------------------------------------------------------------------------

% Check if marker should be inserted by peeking ahead
% This is called when returning to a shallower nesting level
\cs_new_protected:Npn \__contmark_check_and_insert:
  {
    \bool_if:NT \g_contmark_should_check_bool
      {
        \bool_gset_false:N \g_contmark_should_check_bool
        \peek_meaning_ignore_spaces:NTF \begin
          {
            % Next command is \begin (likely another asection) - don't insert marker
          }
          {
            \peek_meaning_ignore_spaces:NTF \end
              {
                % Parent section closing - don't insert marker
              }
              {
                % Actual content follows - insert marker
                \continuationmarker
              }
          }
      }
  }

% Section dispatch logic
% ----------------------------------------------------------------------------
\cs_new_protected:Npn \__contmark_dispatch_section:nn #1#2
  {
    \int_case:nnF {#1}
      {
        {1}{\section{#2}}
        {2}{\subsection{#2}}
        {3}{\subsubsection{#2}}
        {4}{\paragraph{#2}}
        {5}{\subparagraph{#2}}
        {6}{\subsubparagraph{#2}}
      }
      {\paragraph{#2}}
  }

% Nesting management
% ----------------------------------------------------------------------------
\cs_new_protected:Npn \__contmark_begin_section:n #1
  {
    \stepcounter{asectiondepth}
    \int_compare:nNnT { \value{asectiondepth} } > { \value{maxasectiondepth} }
      {
        \setcounter{maxasectiondepth}{\value{asectiondepth}}
      }
    \__contmark_dispatch_section:nn { \value{asectiondepth} } {#1}
  }

\cs_new_protected:Npn \__contmark_end_section:
  {
    \addtocounter{asectiondepth}{-1}
    \int_compare:nNnT { \value{asectiondepth} } < { \value{maxasectiondepth} }
      {
        \int_compare:nNnT { \value{asectiondepth} } > { 0 }
          {
            \bool_gset_true:N \g_contmark_should_check_bool
            \setcounter{maxasectiondepth}{\value{asectiondepth}}
          }
      }
    \__contmark_check_and_insert:
  }

\newenvironment{asection}[1]
  {
    \__contmark_begin_section:n {#1}
  }
  {
    \__contmark_end_section:
  }
\ExplSyntaxOff
